\documentclass[12pt]{article}
\usepackage[utf8]{inputenc}
\usepackage[T1]{fontenc}
\usepackage[margin=1in]{geometry}
\usepackage{setspace}
\usepackage{xcolor}
\usepackage[colorlinks=true, urlcolor=blue!60!black, linkcolor=blue!60!black]{hyperref}
\setlength{\parindent}{2em}
\setlength{\parskip}{1.25em}
\renewcommand{\baselinestretch}{1.0}
\begin{document}
\begin{center}
{\LARGE \textit{303 Creative LLC v.\ Elenis}, 600 U.S.\ 570 (2023)\par}
\vspace{0.8em}
{\large Case Brief\par}
\end{center}
\vspace{1em}

\section{Memory Jogger}
The First Amendment's Free Speech Clause prohibits a state from using its public accommodations law to compel a commercial website designer to create customized, expressive wedding websites conveying messages that conflict with her beliefs (\href{https://supreme.justia.com/cases/federal/us/600/21-476/}{Justia}; \href{https://www.supremecourt.gov/opinions/22pdf/21-476_c185.pdf}{Official Slip Opinion}).

\section{Detailed Case Facts}
Lorie Smith, the founder and sole member-owner of 303 Creative LLC, operates a Colorado graphic and website design business (\href{https://supreme.justia.com/cases/federal/us/600/21-476/}{Justia}). Wishing to expand into wedding websites featuring original text, graphics, and videos telling each couple's love story, Smith sought to serve only opposite-sex weddings because of her sincerely held Christian belief that marriage is reserved to unions between one man and one woman (\href{https://supreme.justia.com/cases/federal/us/600/21-476/}{Justia}). She also wanted to post a statement on her website explaining her religious beliefs and her resulting refusal to create websites for same-sex weddings (\href{https://www.everycrsreport.com/reports/LSB11000.html}{CRS Report LSB11000}). At the time she sued, 303 Creative had not yet designed or sold any wedding websites (\href{https://supreme.justia.com/cases/federal/us/600/21-476/}{Justia}).

\subsection*{Statutory and Constitutional Framework}
\begin{itemize}
  \item \textbf{Colorado Anti-Discrimination Act (CADA):} Originally enacted in 1885 and expanded in 2008 to cover sexual orientation, CADA defines a ``place of public accommodation'' broadly to encompass ``any place of business engaged in any sales to the public and any place offering services \ldots{} to the public,'' excluding places principally used for religious purposes (Colo.\ Rev.\ Stat.\ \S~24-34-601(1)) (\href{https://supreme.justia.com/cases/federal/us/600/21-476/}{Justia}).
  \item \textbf{Accommodation Clause:} Prohibits a public accommodation from directly or indirectly refusing, withholding, or denying ``the full and equal enjoyment'' of its goods or services to any person based on race, creed, disability, sex, sexual orientation, or other protected traits (Colo.\ Rev.\ Stat.\ \S~24-34-601(2)(a)) (\href{https://supreme.justia.com/cases/federal/us/600/21-476/}{Justia}).
  \item \textbf{Communication Clause:} Prohibits a public accommodation from publishing any communication indicating that goods or services will be refused or that an individual's patronage is ``unwelcome, objectionable, unacceptable, or undesirable'' based on a protected trait (Colo.\ Rev.\ Stat.\ \S~24-34-601(2)(a)) (\href{https://supreme.justia.com/cases/federal/us/600/21-476/}{Justia}).
  \item \textbf{Enforcement and Remedies:} Either state officials or private citizens may bring enforcement actions (Colo.\ Rev.\ Stat.\ \S\S~24-34-306, 24-34-602(1)), carrying penalties including fines up to \$500 per violation, cease-and-desist orders, and affirmative remedies such as mandatory educational programs and compliance reporting (\S\S~24-34-306(9), 24-34-602(1)(a), 24-34-605) (\href{https://supreme.justia.com/cases/federal/us/600/21-476/}{Justia}).
  \item \textbf{Constitutional Provision:} Free Speech Clause of the First Amendment (incorporated against the states via the Fourteenth Amendment).
\end{itemize}

\subsection*{Key Stipulated Facts}
To facilitate summary judgment, Smith and Colorado entered sweeping stipulations (\href{https://supreme.justia.com/cases/federal/us/600/21-476/}{Justia}):
\begin{itemize}
  \item Smith is willing to work with clients of any sexual orientation and will gladly create custom graphics and websites for LGBTQ+ clients so long as the content does not violate her beliefs.
  \item She will not produce content contradicting ``biblical truth'' regardless of who requests it.
  \item Every website 303 Creative designs is an ``original, customized'' and ``expressive'' creation designed in close collaboration with the couple to ``celebrate and promote'' her view of marriage, displaying 303 Creative's name so viewers know it is her original artwork.
  \item Numerous other companies in Colorado and nationwide offer custom website design services.
\end{itemize}

\subsection*{Historical and Litigation Context}
The suit followed Colorado's enforcement of CADA against baker Jack Phillips in \textit{Masterpiece Cakeshop, Ltd.\ v.\ Colorado Civil Rights Commission}, 584 U.S.\ 617 (2018), which the Supreme Court had resolved narrowly on Free Exercise grounds (religious hostility by state commissioners) without deciding the underlying Free Speech question (\href{https://www.everycrsreport.com/reports/LSB11000.html}{CRS Report LSB11000}).

\section{Procedural History}
\begin{itemize}
  \item \textbf{U.S.\ District Court for the District of Colorado (2016--2019):} Smith filed a pre-enforcement suit seeking declaratory and injunctive relief under the Free Speech and Free Exercise Clauses (\href{https://supreme.justia.com/cases/federal/us/600/21-476/}{Justia}). The district court initially dismissed her Accommodation Clause challenge for lack of standing and later granted summary judgment to Colorado on her Communication Clause challenge (\textit{303 Creative LLC v.\ Elenis}, 405 F.\ Supp.\ 3d 907 (D.\ Colo.\ 2019)), reasoning that the Communication Clause validly barred advertising conduct made unlawful by the Accommodation Clause (\href{https://www.justice.gov/d9/case-documents/attachments/2022/08/19/303_creative_llc_v._elenis_no._21-476_merits_08-19-22.pdf}{DOJ Merits Brief}; \href{https://www.journals.uchicago.edu/doi/pdf/10.1086/728867}{Robert Post, \textit{The Supreme Court Review}}).
  \item \textbf{U.S.\ Court of Appeals for the Tenth Circuit (2021):} A divided panel affirmed (\textit{303 Creative LLC v.\ Elenis}, 6 F.4th 1160 (10th Cir.\ 2021)) (\href{https://supreme.justia.com/cases/federal/us/600/21-476/}{Justia}). Reviewing standing de novo, the Tenth Circuit held Smith established a ``credible threat'' of enforcement under \textit{Susan B.\ Anthony List v.\ Driehaus}, 573 U.S.\ 149 (2014), given Colorado's enforcement history (including \textit{Masterpiece Cakeshop}), private complaint mechanisms, and refusal to disavow enforcement (\href{https://supreme.justia.com/cases/federal/us/600/21-476/}{Justia}). On the merits, the panel agreed Smith's wedding websites were ``pure speech'' triggering strict scrutiny, but held 2--1 that CADA survived strict scrutiny because Colorado had a compelling interest in equal access and Smith's custom services were ``unique'' and ``unavailable elsewhere'' (\href{https://supreme.justia.com/cases/federal/us/600/21-476/}{Justia}). Chief Judge Tymkovich dissented, warning that the panel's logic would let the government regulate the messages of all artists (\href{https://supreme.justia.com/cases/federal/us/600/21-476/}{Justia}).
  \item \textbf{U.S.\ Supreme Court (2022--2023):} The Court granted certiorari limited solely to the Free Speech Clause question: whether applying a public-accommodation law to compel an artist to speak or stay silent violates the Free Speech Clause (declining to take up the Free Exercise question) (\href{https://supreme.justia.com/cases/federal/us/600/21-476/}{Justia}). Neither party contested Article III standing before the Supreme Court (\href{https://supreme.justia.com/cases/federal/us/600/21-476/}{Justia}). On June 30, 2023, the Court reversed 6--3 (\href{https://supreme.justia.com/cases/federal/us/600/21-476/}{Justia}).
  \item \textbf{Remand:} On remand, the district court entered a permanent injunction barring Colorado from enforcing CADA's Accommodation and Communication Clauses to compel Smith to create same-sex wedding websites or stop her from posting her explanatory statement, and awarded attorney's fees under 42 U.S.C.\ \S~1988(b).
\end{itemize}

\section{Judicial Votes}
\begin{itemize}
  \item \textbf{Majority (6):} Justice Gorsuch (author), joined by Chief Justice Roberts and Justices Thomas, Alito, Kavanaugh, and Barrett (\href{https://supreme.justia.com/cases/federal/us/600/21-476/}{Justia}).
  \item \textbf{Concurrence(s):} None (\href{https://supreme.justia.com/cases/federal/us/600/21-476/}{Justia}).
  \item \textbf{Dissent (3):} Justice Sotomayor (author), joined by Justices Kagan and Jackson (\href{https://supreme.justia.com/cases/federal/us/600/21-476/}{Justia}).
\end{itemize}

\section{Holding}
\textbf{Judgment:} Reversed (6--3) (\href{https://supreme.justia.com/cases/federal/us/600/21-476/}{Justia}).

\textbf{Rule Announced:} Under the Free Speech Clause of the First Amendment, a state may not apply its public accommodations law to compel a business that creates customized, expressive ``pure speech'' (here, original wedding websites) to design expressive works conveying messages with which the creator disagrees---even when the creator offers those services for profit to the general public through a corporate entity (\href{https://supreme.justia.com/cases/federal/us/600/21-476/}{Justia}).

\textbf{Holding vs.\ Dicta:} The Court's holding is tethered to the parties' stipulations that Smith's websites are customized, original ``pure speech'' reflecting her own voice and that she serves LGBTQ+ clients in non-conflicting contexts; the Court expressly declined to announce a general test for determining which other commercial services (e.g., photographers, bakers, stationers) qualify as ``expressive activity,'' noting that such boundary questions are not presented by the stipulated record (\href{https://supreme.justia.com/cases/federal/us/600/21-476/}{Justia}). Because Colorado conceded that its Communication Clause authority rose or fell with the Accommodation Clause, the Court focused its holding on the Accommodation Clause (\href{https://supreme.justia.com/cases/federal/us/600/21-476/}{Justia}).

\section{Analysis of Opinions}

\subsection{Majority Opinion (Justice Gorsuch)}
Justice Gorsuch structures the opinion around a categorical anti-compulsion principle grounded in the Free Speech Clause (\href{https://supreme.justia.com/cases/federal/us/600/21-476/}{Justia}):
\begin{itemize}
  \item \textbf{Foundational Compelled-Speech Doctrine:} Opening with the principle that the First Amendment safeguards an ``uninhibited marketplace of ideas'' (\textit{McCullen v.\ Coakley}, 573 U.S.\ 464 (2014)) and bars the state from prescribing ``what shall be orthodox'' (\textit{West Virginia State Board of Education v.\ Barnette}, 319 U.S.\ 624 (1943)), Gorsuch relies heavily on a trilogy of cases---\textit{Barnette}, \textit{Hurley v.\ Irish-American Gay, Lesbian and Bisexual Group of Boston, Inc.}, 515 U.S.\ 557 (1995), and \textit{Boy Scouts of America v.\ Dale}, 530 U.S.\ 640 (2000)---to establish that the government may neither force a person to speak a state-preferred message nor use a public accommodations law to alter a speaker's expressive content (\href{https://supreme.justia.com/cases/federal/us/600/21-476/}{Justia}).
  \item \textbf{Application to Stipulated Facts (``Pure Speech'' and Attribution):} Agreeing with the Tenth Circuit, Gorsuch finds that Smith's planned websites are ``pure speech'' conveyed over the internet (\textit{Reno v.\ ACLU}, 521 U.S.\ 844 (1997)) and constitute \emph{her} speech because she vets projects, collaborates with couples, and authors original words and artwork under her company's name; combining her voice with the couple's does not forfeit First Amendment protection (\textit{Hurley}) (\href{https://supreme.justia.com/cases/federal/us/600/21-476/}{Justia}).
  \item \textbf{Rejection of the Tenth Circuit's ``Unique Services'' Rationale:} While acknowledging that states have a ``compelling interest'' in eliminating discrimination in public accommodations (\textit{Roberts v.\ United States Jaycees}, 468 U.S.\ 609 (1984); \textit{Heart of Atlanta Motel, Inc.\ v.\ United States}, 379 U.S.\ 241 (1964)) and may protect gay persons in ordinary commercial transactions (\textit{Masterpiece Cakeshop}), Gorsuch emphasizes that public accommodations laws have expanded far beyond common-law common carriers and innkeepers and cannot override the Constitution (U.S.\ Const.\ art.\ VI, cl.\ 2) (\href{https://supreme.justia.com/cases/federal/us/600/21-476/}{Justia}). Treating an artist's ``unique'' talent as a justification for state coercion would invert the First Amendment by making the finest artists easiest to conscript (\href{https://supreme.justia.com/cases/federal/us/600/21-476/}{Justia}).
  \item \textbf{Rejection of Colorado's Alternative Arguments:}
    \begin{itemize}
      \item \emph{Commercial / For-Profit Status:} Accepting payment or operating as an LLC does not strip First Amendment protection (citing \textit{Joseph Burstyn, Inc.\ v.\ Wilson}, 343 U.S.\ 495 (1952); \textit{Simon \& Schuster, Inc.\ v.\ Members of N.Y.\ State Crime Victims Bd.}, 502 U.S.\ 105 (1991)) (\href{https://supreme.justia.com/cases/federal/us/600/21-476/}{Justia}).
      \item \emph{Status vs.\ Message:} Pointing to the stipulations that Smith serves LGBTQ+ clients for other projects, Gorsuch characterizes her objection as message-based rather than status-based, adding that speaker motive does not forfeit First Amendment protection (\textit{Snyder v.\ Phelps}, 562 U.S.\ 443 (2011); \textit{National Socialist Party of America v.\ Skokie}, 432 U.S.\ 43 (1977)) (\href{https://supreme.justia.com/cases/federal/us/600/21-476/}{Justia}).
      \item \emph{Distinguishing FAIR:} Unlike \textit{Rumsfeld v.\ Forum for Academic \& Institutional Rights, Inc.} (\textit{FAIR}), 547 U.S.\ 47 (2006)---where requiring law schools to send logistical recruiter emails was an ``incidental'' burden on conduct---forcing Smith to create custom celebratory wedding websites directly compels speech on a matter of political and religious significance (\href{https://supreme.justia.com/cases/federal/us/600/21-476/}{Justia}).
    \end{itemize}
  \item \textbf{Slippery-Slope Counter-Examples:} Echoing Chief Judge Tymkovich, Gorsuch warns that Colorado's rule would let states force a Muslim film director to make a Zionist film, an atheist muralist to celebrate Evangelical zeal, or a gay website designer to build a site opposing same-sex marriage (\href{https://supreme.justia.com/cases/federal/us/600/21-476/}{Justia}).
\end{itemize}

\subsection{Dissenting Opinion (Justice Sotomayor)}
Justice Sotomayor frames the case as an unprecedented departure from the historical social contract of public accommodations law and settled conduct-regulation doctrine (\href{https://supreme.justia.com/cases/federal/us/600/21-476/}{Justia}):
\begin{itemize}
  \item \textbf{History and Dual Purposes of Public Accommodations Laws:} Tracing the common-law duty of businesses that ``hold themselves out'' to serve the public (\textit{Lane v.\ Cotton}, 88 Eng.\ Rep.\ 1458 (K.B.\ 1701)) through the Civil Rights Act of 1964, sex-equality laws, the ADA, and LGBT protections, Sotomayor argues that public accommodations laws serve two compelling interests: equal access to the market and equal dignity against stigmatic harm (\textit{Heart of Atlanta Motel}; \textit{Roberts}) (\href{https://supreme.justia.com/cases/federal/us/600/21-476/}{Justia}). She rejects Gorsuch's claim that common-law duties rested on monopoly power (\href{https://supreme.justia.com/cases/federal/us/600/21-476/}{Justia}).
  \item \textbf{Precedent Rejecting Constitutional Rights to Discriminate:} Sotomayor cites \textit{Heart of Atlanta Motel}, \textit{Katzenbach v.\ McClung}, 379 U.S.\ 294 (1964) (rejecting Ollie's Barbecue's ``limited menu'' offer of takeout to Black patrons), \textit{Newman v.\ Piggie Park Enterprises, Inc.}, 390 U.S.\ 400 (1968), \textit{Runyon v.\ McCrary}, 427 U.S.\ 160 (1976) (commercially operated schools had no First Amendment right to exclude Black students), \textit{Hishon v.\ King \& Spalding}, 467 U.S.\ 69 (1984) (law firm had no First Amendment right to deny partnership to a woman despite expressive legal advocacy), and \textit{Roberts} to show the Court has consistently rejected First Amendment exemptions for commercial refusals of service (\href{https://supreme.justia.com/cases/federal/us/600/21-476/}{Justia}).
  \item \textbf{Conduct Regulation and \textit{O'Brien} / \textit{FAIR} Framework:} CADA regulates the commercial \emph{act} of refusing service based on status, not the content of speech. Under \textit{United States v.\ O'Brien}, 391 U.S.\ 367 (1968), \textit{Sorrell v.\ IMS Health Inc.}, 564 U.S.\ 552 (2011), and \textit{FAIR}, any burden on speech is ``incidental'' because speech is compelled only ``if, and to the extent,'' the business chooses to sell such services to other members of the public (\href{https://supreme.justia.com/cases/federal/us/600/21-476/}{Justia}). Under CADA, Smith remains free to include biblical quotes or refuse specific slogans (``Love is Love'') for all customers; she simply cannot refuse the same wedding website to a same-sex couple solely because of who is marrying (\href{https://supreme.justia.com/cases/federal/us/600/21-476/}{Justia}).
  \item \textbf{Distinguishing \textit{Hurley}, \textit{Dale}, and \textit{Barnette}:} \textit{Hurley} and \textit{Dale} involved ``peculiar'' applications of public accommodations laws to private, nonprofit expressive associations, not ``clearly commercial entities'' selling services to the public; equating a neutral equal-access rule with \textit{Barnette}'s forced flag salute ``trivializes the freedom protected in \textit{Barnette}'' (\textit{FAIR}) (\href{https://supreme.justia.com/cases/federal/us/600/21-476/}{Justia}).
  \item \textbf{Rejection of Status-vs.-Message Distinction and Future Risks:} Sotomayor argues that defining a service as ``opposite-sex wedding websites'' collapses status into message---just as a hotel could not recast its service as ``whites-only lodgings'' or a photo studio sell ``passport photos for white people.'' Because the holding rests on the Free Speech Clause rather than Free Exercise, its logic is not limited to religious objections or sexual orientation, threatening refusals of wedding websites for interracial couples, birth announcements for disabled parents, or headstone inscriptions for same-sex partners (\href{https://supreme.justia.com/cases/federal/us/600/21-476/}{Justia}).
\end{itemize}

\section{Examples: Future Applications}

\subsection*{Same-Side (Protected Under \textit{303 Creative})}
\begin{enumerate}
  \item \textbf{Bespoke Wedding Photographer:} A commercial wedding photographer who shoots, curates, and edits custom visual narratives of weddings refuses to photograph a same-sex wedding on religious grounds while willingly taking non-wedding portraits for LGBTQ+ clients. Because custom photography is original visual expression communicating a celebratory message attributable to the artist, compelling the photographer to shoot and edit the wedding compels ``pure speech'' under \textit{303 Creative}. After \textit{303 Creative}, the Second Circuit held that a wedding photographer raising this claim had plausibly stated a free speech claim at the pleading stage, while declining to enter a preliminary injunction and remanding for further proceedings, in \textit{Emilee Carpenter, LLC v.\ James}, 107 F.4th 92 (2d Cir.\ 2024) (\href{https://uclawreview.org/2026/10/16/when-does-free-expression-end-and-discrimination-begin-an-analysis-of-how-lower-courts-have-applied-the-303-creative-llc-v-elenis-decision/}{UC Law Review Blog}).
  \item \textbf{Freelance Speechwriter / Ghostwriter Open to the Public:} A solo commercial speechwriter who advertises custom wedding toasts and corporate keynotes to the general public refuses a commission to write a speech celebrating an evangelical religious revival because she is an atheist. Even though she operates for profit and holds herself out to the public, writing original prose conveying another's religious message is quintessential pure speech that a public accommodations law cannot compel.
\end{enumerate}

\subsection*{Opposite-Side (Not Protected Under \textit{303 Creative})}
\begin{enumerate}
  \setcounter{enumi}{2}
  \item \textbf{Wedding Reception Venue / Catering Hall:} A hotel ballroom or rustic barn that rents event space and provides standard banquet catering to the public refuses to host a same-sex wedding reception. Renting physical space and serving food are ordinary commercial goods and services---not ``pure speech'' conveying the owner's own original message---so CADA and similar public accommodations laws validly regulate this commercial conduct without violating the First Amendment (\href{https://supreme.justia.com/cases/federal/us/600/21-476/}{Justia}).
  \item \textbf{Standardized Print Shop / Off-the-Shelf Invitation Retailer:} A commercial print shop that lets customers upload their own finished PDF wedding invitations for automated printing, or sells pre-printed template cards where the shop merely swaps in names and dates without original design input, refuses to print invitations for a same-sex couple. Because the shop is selling a standardized commercial product rather than authoring customized, original artwork attributable to the business, the refusal is regulable status-based conduct rather than compelled speech.
\end{enumerate}

\subsection*{Fence-Sitter (Genuinely Unclear Outcome)}
\begin{enumerate}
  \setcounter{enumi}{4}
  \item \textbf{Custom Wedding Cake Baker or Bespoke Wedding Florist:} A baker who designs elaborate, hand-sculpted custom wedding cakes (or a florist who creates bespoke floral installations) refuses to create a custom wedding cake for a same-sex couple, though willing to sell them pre-made cakes or birthday cakes. Unlike websites or written prose, cakes and floral arrangements do not automatically qualify as ``pure speech,'' and \textit{303 Creative} expressly left open how courts should determine when non-verbal, partly functional crafts cross the line into ``expressive activity'' attributed to the creator (\href{https://supreme.justia.com/cases/federal/us/600/21-476/}{Justia}; \href{https://www.everycrsreport.com/reports/LSB11000.html}{CRS Report LSB11000}). Lower courts must wrestle with whether a reasonable observer perceives a custom cake or bouquet as the vendor's own message celebrating a particular conception of marriage or merely as an artisanal commercial product.
\end{enumerate}

\section{Critique}
\begin{itemize}
  \item \textbf{Doctrinal Ambiguity and the Missing Scrutiny Framework:} Scholars note that Gorsuch's majority opinion never clearly specifies what tier of scrutiny it applies---effectively bypassing strict scrutiny in favor of a categorical bar against compelled ``pure speech'' (\href{https://www.everycrsreport.com/reports/LSB11000.html}{CRS Report LSB11000}; \href{https://www.journals.uchicago.edu/doi/pdf/10.1086/728867}{Robert Post, \textit{The Supreme Court Review}}). Robert Post criticizes the opinion's ``impossible abstraction,'' arguing that by treating ``pure speech'' as an automatic trump without balancing competing constitutional values or defining the boundary between commercial conduct and public discourse, the Court creates doctrinal chaos for lower courts (\href{https://www.journals.uchicago.edu/doi/pdf/10.1086/728867}{Robert Post}). Similarly, commentators observe that \textit{303 Creative} fails to adopt a workable ``reasonable-audience'' or attribution test to distinguish genuine compelled endorsement from routine commercial service (\href{https://uclawreview.org/2026/10/16/when-does-free-expression-end-and-discrimination-begin-an-analysis-of-how-lower-courts-have-applied-the-303-creative-llc-v-elenis-decision/}{UC Law Review Blog}).
  \item \textbf{Pre-Enforcement Posture, Stipulations, and the ``Stewart'' Controversy:} Because Smith had never designed a wedding website, the entire case turned on broad, abstract stipulations that allowed the majority and dissent to ``talk past each other''---one seeing a mandate to endorse a contested ideology, the other seeing a refusal to sell an identical service to gay couples (\href{https://scholarship.law.georgetown.edu/cgi/viewcontent.cgi?article=3596&context=facpub}{David Cole, \textit{Yale Law Journal Forum}}). Just before the decision issued, \href{https://newrepublic.com/article/173987/mysterious-case-fake-gay-marriage-website-real-straight-man-supreme-court}{\textit{The New Republic} reported} that a purported wedding-website inquiry from a man named ``Stewart'' cited in lower-court filings belonged to a married straight man who denied ever contacting 303 Creative (\href{https://www.cbsnews.com/news/supreme-court-303-creative-case-stewart-wedding-website-lgbtq-free-speech/}{CBS News}). While defenders noted that the district court had discounted the inquiry and that neither the Tenth Circuit nor the Supreme Court relied on it to find pre-enforcement standing (\href{https://reason.com/volokh/2023/07/18/more-on-standing-in-the-303-creative-case/}{Volokh Conspiracy / Reason}; \href{https://adflegal.org/article/responding-media-smear-303-creative/}{ADF}), critics argued the contrived pre-enforcement record deprived the Court of concrete facts needed to test whether Colorado would actually compel ideological speech (\href{https://adamunikowsky.substack.com/p/contrived-cases-make-bad-law}{Adam Unikowsky}; \href{https://scholarship.law.georgetown.edu/cgi/viewcontent.cgi?article=3596&context=facpub}{David Cole}).
  \item \textbf{Progressive / Egalitarian Critiques:} Progressive scholars argue the Court weaponized First Amendment doctrine to carve a deregulatory hole in civil rights law (\href{https://www.californialawreview.org/print/equal-right-to-exclude}{\textit{California Law Review}}) and misunderstood how commercial antidiscrimination laws operate by targeting transaction terms rather than expression (\href{https://digitalcommons.law.seattleu.edu/cgi/viewcontent.cgi?article=2909&context=sulr}{Helen Norton, \textit{Seattle University Law Review}}). Michael Dorf also highlights an asymmetry of class and ownership: the First Amendment shield protects \emph{business owners} who object to a client's message, whereas \emph{employees} of an expressive business receive no First Amendment protection if fired by a private boss for refusing to create a message they oppose (\href{https://www.dorfonlaw.org/2022/12/a-selective-preview-of-scotus-oral.html}{Dorf on Law}). At the same time, some civil libertarians (including ACLU Legal Director David Cole) contend that \textit{303 Creative} is properly read narrowly---protecting only truly customized, original expression---and should leave the vast majority of public accommodations laws intact (\href{https://scholarship.law.georgetown.edu/cgi/viewcontent.cgi?article=3596&context=facpub}{David Cole}).
  \item \textbf{Originalist / Textualist Perspectives:} Originalist commentators diverge on the opinion's methodology. Some argue the majority largely bypassed original public meaning in favor of 20th-century compelled-speech precedents (\textit{Barnette}, \textit{Hurley}, \textit{Dale}), leaving unresolved whether the Founding-era First Amendment barred states from regulating commercial trades (\href{https://houstonhealthlaw.scholasticahq.com/article/168216-mahmoud-skrmetti-and-303-creative-ignoring-original-meaning-rewriting-precedent-and-discounting-harm-to-lgbtq-people/attachment/361335.pdf}{\textit{Houston Journal of Health Law \& Policy}}). Conversely, originalist scholar Robert G.\ Natelson argues that an originalist analysis would have made the case \emph{easier} for Smith: in the Founding era, ``freedom of the press'' specifically protected for-profit publication through a medium (like newspapers and printers), eliminating the modern doctrinal friction over whether commercial enterprises enjoy full First Amendment protection (\href{https://fedsoc.org/commentary/fedsoc-blog/with-an-originalist-understanding-of-the-first-amendment-the-303-creative-case-would-have-been-much-easier}{Federalist Society Blog}).
  \item \textbf{Competing Core Values:} Ultimately, the 6--3 divide reflects a clash of foundational premises. If one prioritizes individual intellectual autonomy and distrusts state power to police orthodoxy in any expressive medium, Gorsuch's rule follows naturally (\href{https://www.thepublicdiscourse.com/2023/07/89894/}{\textit{The Public Discourse}}). If one prioritizes the anti-caste social contract---the premise that entering the public marketplace carries a reciprocal duty to serve all citizens with equal dignity---Sotomayor's conduct-regulation framework is more compelling (\href{https://supreme.justia.com/cases/federal/us/600/21-476/}{Justia}).
\end{itemize}

\section{Key Quotations}
All quotations below are verified directly against the published Supreme Court opinions (\href{https://supreme.justia.com/cases/federal/us/600/21-476/}{Justia}):

\begin{quote}
``When a state public accommodations law and the Constitution collide, there can be no question which must prevail.''
\end{quote}
\noindent --- Justice Gorsuch, majority opinion (Supremacy of the Constitution over Public Accommodations Laws).

\begin{quote}
``Were the rule otherwise, the better the artist, the finer the writer, the more unique his talent, the more easily his voice could be conscripted to disseminate the government's preferred messages. That would not respect the First Amendment; more nearly, it would spell its demise.''
\end{quote}
\noindent --- Justice Gorsuch, majority opinion (Rejecting the Tenth Circuit's ``Unique Services'' Theory).

\begin{quote}
``But tolerance, not coercion, is our Nation's answer. The First Amendment envisions the United States as a rich and complex place where all persons are free to think and speak as they wish, not as the government demands.''
\end{quote}
\noindent --- Justice Gorsuch, majority opinion (Core Free Speech Vision).

\begin{quote}
``Today, the Court, for the first time in its history, grants a business open to the public a constitutional right to refuse to serve members of a protected class.''
\end{quote}
\noindent --- Justice Sotomayor, dissenting (Unprecedented Exemption for Public Businesses).

\begin{quote}
``As I will explain, the law in question targets conduct, not speech, for regulation, and the act of discrimination has never constituted protected expression under the First Amendment. Our Constitution contains no right to refuse service to a disfavored group.''
\end{quote}
\noindent --- Justice Sotomayor, dissenting (Conduct vs.\ Speech).

\end{document}
